\chapter*{THÔNG TIN BÀI BÁO GỐC}
\addcontentsline{toc}{chapter}{THÔNG TIN BÀI BÁO GỐC}
\noindent\textbf{Tên bài báo:} Improving stroke risk prediction by integrating XGBoost, optimized principal component analysis, and explainable artificial intelligence.\\[4pt]
\textbf{Tác giả:} Lesia Mochurad, Viktoriia Babii, Yuliia Boliubash và Yulianna Mochurad.\\[4pt]
\textbf{Tạp chí:} BMC Medical Informatics and Decision Making (2025) 25:63.\\[4pt]
\textbf{DOI:} \url{https://doi.org/10.1186/s12911-025-02894-z}.\\[8pt]
\textbf{Nguồn dữ liệu:}
\begin{itemize}
  \item \url{https://www.kaggle.com/fedesoriano/stroke-prediction-dataset}
  \item \url{https://www.kaggle.com/datasets/pranavp1999/stroke-prediction-health-care-synthetic-dataset}
\end{itemize}

\chapter*{GIỚI THIỆU BÁO CÁO}
\addcontentsline{toc}{chapter}{GIỚI THIỆU BÁO CÁO}
Báo cáo này được thực hiện trong khuôn khổ học phần Khai phá dữ liệu, với
mục tiêu tìm hiểu và trình bày lại một công trình nghiên cứu ứng dụng học máy, chúng tôi sẽ thực hiện tái hiện bài nghiên cứu và để xuất các điểm có thể tối ưu hoặc cái tiến của bài báo gốc. Lưu ý rằng từ Chương 1 đến hết Chương 4, nội dung báo cáo chủ yếu là phần dịch và
trình bày lại bài báo gốc. Do đó, đại từ \textit{``chúng tôi''} trong
các chương này được giữ theo nguyên bản và dùng để chỉ nhóm tác giả của bài báo gốc. Từ Chương 5 trở đi, đại từ ``chúng tôi'' dùng để chỉ nhóm học viên thực hiện báo cáo này.

\medskip

\noindent
\textbf{Lý do lựa chọn bài báo:} đây là một nghiên cứu tích hợp đồng thời ba
kỹ thuật đã được giới thiệu trong học phần: (1) thuật toán học máy có giám sát mạnh
(XGBoost), (2) kỹ thuật giảm chiều dữ liệu (PCA), và (3) trí tuệ nhân tạo giải
thích được (XAI, cụ thể là SHAP). Việc bài báo kết hợp cả ba khía cạnh này
trong một quy trình xử lý dữ liệu y tế thực tế giúp minh họa rõ mối liên hệ
giữa lý thuyết khai phá dữ liệu và một ứng dụng có ý nghĩa thực tiễn cao.

\medskip

\noindent
\textbf{Mục tiêu của báo cáo:}

\begin{itemize}
  \item Trình bày lại có hệ thống nội dung bài báo gốc bằng tiếng Việt.
  \item Tự xây dựng và chạy lại pipeline thực nghiệm nhằm tái lập các kết quả
  mà tác giả công bố trên hai bộ dữ liệu.
  \item Đối chiếu có kiểm soát giữa kết quả tác giả và kết quả tái hiện của nhóm,
  từ đó xác định các thành phần của pipeline cần hiệu chỉnh để tiến gần kết quả gốc.
  \item Sau khi thiết lập được baseline tái hiện phù hợp, khảo sát các hướng cải tiến
  và đánh giá liệu chúng có tạo ra lợi ích thực nghiệm so với phương pháp của tác giả hay không.
\end{itemize}